\section*{はじめに}
昨今はデータサイエンス，情報科学の領域が非常に隆盛で，コンピュータを使ってデータを分析し，経済の動向や購買行動などの予測に用いられることが広く行われている。

人の行動や考え方をどのようにデータにするかについては，当然ながら心理学には一日の長がある。
また，人が頭の中でどのような考え方のプロセスをたどるのか，それをどのように検証するのかについても，心理学はその短い歴史の中で徹底的にその技法を洗練させてきた。
このような根源的なレベルでの理論や方法論は時代が変わっても色褪せることなく，また今後ますます必要とされてくる時代になっている。

本講ではデータ解析の応用段階として，より実践的なテーマを扱う。すなわち，\textbf{心理尺度が作られる理論的背景}と，\textbf{データの背後のメカニズムを解析する方法}を知ることである。

心理学研究法の1つとして，調査研究がある。紙とペンで回答を集めた時，回答者がある反応カテゴリにまるをつけたことが，どうして数値処理の対象になるのか。そこには数字を割り振るルールとしての「尺度化」の手続きがある。残念ながら応用的側面が発展しすぎたため，回答に数字を割り振る原理について語られることが少なくなってきてしまい，それに対する反動からか，近年改めてこの根本原理についての理解と解説が求められている。この講義では，尺度化の原理や目に見えない潜在変数を想定して分析するとはどういうことかについて，理論と演習を交えながら習得することを目指す。この理論的側面を考えるためには，どうしても線形代数・行列計算の知識が必要になってくる。線形代数については特別な事前知識は不要で，定義から改めて解説するので安心してもらいたい。

心理学研究法のまた1つの大きな柱として，実験的研究がある。実験的研究はその手続きが厳格に準備されていることで，結果は群間の平均値差を推定すれば十分である，とされてきた。心理学の基礎領域では，そのため，標本の平均値から母集団の平均値を推測する方法が主なテーマになる。しかしこの方法は逆に，結果を群間の平均値差に帰着させるための実験計画を必要とする。結果に方法が規定されているのである。本来考えたかったことは群間の差だけではなく，どのようにして心理的なメカニズムからデータが生成されているか，という問いであったことを思い出そう。そして群間の平均値に拘ることなく，心理的なメカニズムを数式的に表現することで分析する\textbf{数理モデリング}というアプローチがある。このモデリングの基礎的な知識，方法論の習得を目指す。

\subsection*{授業のテーマ}
データから意味のある情報を取り出すための，さまざまな分析法を習熟するにあたって，その背後にあって語られることのない「発想」の観点から理解する。
数値だけに振り回される状態から脱却し，数値を算出する数式に込められた意味について考える視点を持つ。
またこれらに習熟することで，どのような研究対象に対してどのような心理統計的アプローチができるかを，俯瞰的に見れるようになろう。

\subsection*{一年を通じて伝えたいポイント}
\begin{description}
	\item[尺度化とは何か] 心理学で行われるアンケート調査やその後の分析はどういう原理があって「心を測定した」といえるのか。その原理やモデルを理解して利用できるようになる。
	\item[多変量解析から何がわかるのか] 調査研究などで得られた多変量を分析することで何がわかるのか。あるいは何をしてわかったというのか。
	\item[多変量解析の基礎となる数式的原理] 多変量解析の背景にあるのは線形代数という数学であり，線形代数の基礎を学ぶことで多変量解析のメカニズムを統合的に理解できる。
	\item[データ生成メカニズム] データから情報を取り出す受け身的な分析ではなく，データに数字を与えたり，データが生まれてくるメカニズムをリバースエンジニアリングすることで，さらに積極的にデータ解析に立ち向かおう。
	\item[統計環境Rと確率的プログラミング言語Stanによる実践] 統計環境Rと確率的プログラミング言語Stanに習熟することで実際に計算し，確認しながら分析を進めることができる。
\end{description}